\subsubsection{Constructive determination and the function of the dodecaphasic record}
\label{alpha-ampl-dependencias}

The preceding definition makes it possible to specify what it means to determine alpha
from the structure. The domain of the operation contains four records of
known provenance: the three regional sequences and the integral phase
record. Normalization additionally incorporates a boundary condition
compatible with continuation. Once these data are fixed, Euclidean
division determines each remainder and its quotient. Proposition
\ref{prop:alpha-normalizacion} extends this determination to the complete
system. Independence from a target physical value is therefore
a property of the construction order: the value compared at the
end does not participate in selecting any argument of this
map.

This causal independence must be distinguished from algebraic independence
between numbers. The former is established by identifying the domain,
operations, and provenance of the coefficients; the latter would be
a statement about polynomial relations. The result used here is the
former. The construction determines a coordinate through records
produced by APP--TRIT--TPK and subsequently permits verification of the
relations it satisfies. This order makes it possible to use an already
generated coordinate in a later operation without turning it into an external parameter.

The function of the record $K$ is more precise than that of a numerical
correction added to a formula. Its twelve positions originate in the
reversible transformation of the signed record. The equation
$U=\mathcal H_{12}K$ preserves the possibility of returning to that record,
while the transverse differences and total sum allow another
reconstruction. In periodic evaluation, the order of the positions
determines the positional weight of each component. The same set of
components, arranged in a different order, generally produces a different rational
coordinate. The phase origin and orientation therefore form part of the
definition of the number occurring in \eqref{eq:alpha-identidad-completa}.

Alpha normalization composes this information with the regions. Its
local identity distinguishes two successive operations. First, the
oriented balance of blocks is formed; then that balance is represented by
a canonical remainder and a quotient linking adjacent positions. The first
operation preserves the sign of the balance and allows a support to be
extracted subsequently. The second allows prefixes of the real coordinate
to be formed. Both readings remain related because the state preserves the
data before and after normalization. The signed support and the
positional expansion are thus different realizations of an explicit
arithmetic composition.

The telescoping identity expresses the global content of this composition.
The interior quotients cancel in pairs when a prefix is evaluated,
while the endpoint terms remain in the equality. This cancellation is
a property of the weighted sum; it neither removes the quotients from the state nor
allows them to be discarded before the sum is performed. The right boundary of a
window receives the contribution of the continuation. For this reason,
evaluating an isolated window and restricting a continued section
are different operations. The change of the final block from $663$ to
$662$, proved above, provides an explicit realization of this
difference without altering the already stabilized blocks to its left.

In this sense, bidirectionality has a concrete arithmetic
formulation. Generation provides compatible continuations;
normalization transmits the integer information at their boundary from
positions of lower weight towards those of higher weight. The equation preserves the
contribution of both directions through its indices and boundary
conditions. Reversibility in the fixed fibre is verified by reconstructing
$K$ from the other records and the complete quotients. This
inversion certifies the consistency of the original operation; the genealogical
direction remains the one that first produces the record and then
the alpha coordinate.

The second determination has complementary content. In it, the
immediate object is a vacancy balance, and the parameter appears as
its root in an isolating domain. The positional route provides a
canonical normalization of records; the analytic route studies the vanishing
of an operator. Presenting both requires precisely this difference
in function to be retained. Their relationships with the common state must
be expressed through the refinement maps and the comparison of
their complete realizations, as developed in the corresponding
subsection. Agreement at a finite approximation constitutes a
localized check; the identification of the complete constructions
has the domain and scope of its own statement.

Finally, this organization explains alpha's place in the article.
Its equation appears after regional generation and the construction
of $K$, because it uses both results. The electronic construction
uses this coordinate as a factor in its composition. The exceptional
extension receives the records and their supports, in which incidence
information is preserved. These two continuations return to a common
provenance and develop different mathematical functions: evaluating a
central section and constructing a boundary configuration.
Developing each makes it possible to follow the structure through the
equation, rather than restricting the comparison to its scalar value.
